\scanpage{68}
% archive foliation: 67
% lines 1--7, diagonally cancelled calculation
\deleted{
$(p,q,r)$ doit être parmi les \deleted{\unsure{trois}} syst. \unsure{énumérés}

a) $p,2,2$ : $\frac1p+\frac12+\frac12>(1+(1-\frac1s)$\qquad
i.e. $\frac1p+\frac1s>1$ impossible car $\frac1p,\frac1s\le\frac12$

b) $(p,3,2)$ $(p=3,4,5\unsure{)},s\ge p)$
\deleted{$\unsure{\frac1p+\frac13+\frac12}$}
\[
\frac1{\unsure{s|p}}+\frac1p+\underbrace{\frac13+\frac12}_{5/6}>2
\qquad\frac1s+\frac1p>\frac76\quad\text{impossible car }\frac1s,\frac1p\le\frac13
\]
}
% lines 8--11, live calculation
\[
\frac1p+\frac1q>2-\underbrace{\left(\frac1r+\frac1s\right)}_{\substack{\le1\\\text{car }1/r\le1/2\\1/s\le1/2}}\ge1
\quad\text{donc}\quad\frac1p+\frac1q\ge1\quad\text{impossible car }\frac1p,\frac1q\le\frac12
\]
\uline{Impossible} !

% lines 12--14
Car si $\operatorname{Card}I\ge4$, on aura
\[
\frac1{d_1}+\frac1{d_2}+\frac1{d_3}+\frac1{d_4}
\deleted{\unsure{+\sum_{5\le i\le\nu}\frac1{d_i}}}
>\begin{array}{c}\unsure{2}+\sum_{5\le i\le\unsure{\nu}}(1-\frac1{d_i})\\\unsure{\vee}\\2\end{array}
\]
et on vient de voir que c'est impossible !

% lines 15--20, lower draft crossed in large X strokes
\deleted{
\uline{Étude du cas} \deleted{$\unsure{\operatorname{Card}I=2}$}\inserted{$(N,N)$} :
En caractéristique $0$, c'est le cas \unsure{kummérien},
\[G\simeq\mu_N(k)\simeq Z/NZ\qquad\ldots\]
\deleted{\unsure{Plus}} \unsure{gén.}, il en est ainsi si car. $p$ premier à $N$

Si $p\mid N$, on trouve une \inserted{$p$-}extension
\inserted{(de degré $p^r$, \unsure{t.t.} ramifiée et \unsure{inséparable})}
d'une $N'$-extension kummérienne $(N=p^rN',\ (p,N')=1)$.
C'est donc une
% lines 21--26, multiple revisions within the cancelled draft
\uline{extension} \deleted{est} \uline{résoluble},
\deleted{On en peut dire plus,}
\inserted{mais bien mieux : les $p$-\deleted{\unsure{extensions}}sous-groupes
\inserted{$H$} de $\operatorname{GP}(1,k)$ sont abéliens et contenus dans}
\deleted{\unsure{dans un grpe} \illegible{1}. (\unsure{car} \illegible{4})}
\margin{\unsure{conjugués} entre eux, et $H$ isom. à $(Z/pZ)^{\unsure{r|\nu}}$.}
\deleted{\unsure{abélien unipotent}}
sous-groupes \unsure{unip\textsuperscript{ts}} de dim $1$, On trouve donc
% lines 27--30
$G\subset\text{Borel}$, et par suite $G$ produit $\frac12$ direct d'un
groupe cyclique premier à la car., par un $(Z/pZ)^r$.
En fait, on voit aisément (\deleted{\unsure{encore}}
\unsure{puisqu'on est dans le cas} (cyclique premier à car., donc kummérien, qui en est exclu),
on \unsure{tombe en contradiction}.

% margin beside lines 28--32
\margin{En tous les cas, \unsure{on a une} extension abélienne
\unsure{non kummérienne|ou kummérienne}
(d'ordre premier à $p$) \unsure{par un grpe}
$\unsure{(Z/p^nZ)^\nu|(Z/pZ)^\nu}$}

% lines 31--37
\uline{Étude du cas} $(\unsure{N},2,2)$, $(\unsure{[G]}=2N)$.
En car. $\ne2$, \uline{\unsure{l'indice}}\inserted{\unsure{d'ordre} \unsure{vt}}
de ramification $2$ \unsure{en|est} ramification modérée, et faisant une
\deleted{\unsure{d'ordre} \illegible{1}}
extension kummérienne de degré $2$, on \deleted{\unsure{en}} trouve
un $X'/Y'$ qui n'a plus que $2$ pts de ramification \deleted{dans} $Y'$
\diagram{\begin{tikzcd}
X\arrow[d]&X'\arrow[l]\arrow[d,dash]\\
Y&Y'\arrow[l]
\end{tikzcd}}
donc $X'/Y'$ est justiciable de l'étude précédente
mais \unsure{pb} \deleted{\unsure{faut vérifier}} \deleted{\unsure{que si}}
$X'$ de genre $0$ (facile)
}
